% Copyright (c) 2026 Eugene Joseph M. Ragasa.
% SPDX-License-Identifier: Apache-2.0
% Author: Eugene Joseph M. Ragasa
% Formal proof completed: 2026-10-20
\documentclass[11pt]{article}
\usepackage[T1]{fontenc}\usepackage{amsmath,amssymb}\usepackage[margin=1in]{geometry}\usepackage[hidelinks]{hyperref}
\setlength{\emergencystretch}{3em}
\title{Proof Dossier: FCC--Bain Determinant One}\author{Eugene Joseph M. Ragasa}\date{Formal proof completed: 2026-10-20}
\begin{document}\maketitle
\section{Identity and status}
ID \texttt{PKF-PHYSKIT-THM-003}; status \texttt{proved}. Lean declaration:
\par\noindent{\tiny\path{ProjectKoios.Crosswalks.PhysKit.bainToFccMatrix_det}}.
\section{Source and statement}
PhysKit commit \path{7157dcedd0d220673d443c1c025e35a6fd1ffa7a} declares the exact integer matrix
\[
U=\begin{pmatrix}1&1&0\\0&0&1\\0&-1&0\end{pmatrix}
\]
and reports determinant one. The Lean theorem states $\det U=1$ over $\mathbb Z$.
\section{Proof}
Expanding the $3\times3$ determinant gives
$1(0\cdot0-1(-1))-1(0\cdot0-1\cdot0)+0=1$.
Lean uses Mathlib's three-dimensional determinant formula and exact integer normalization.
\section{Boundary and trust}
This proves unimodularity of the encoded constant, not agreement of independently supplied
lattice bases or any physical model. Kernel axioms are \texttt{propext},
\texttt{Classical.choice}, and \texttt{Quot.sound}; no admissions. AI-assisted and
kernel-checked. No independent expert review is recorded.
\end{document}
